\textbf{Cálculo de la Potencia Consumida}

Para el cálculo de la potencia consumida para mover la materia orgánica es importante determinar la velocidad a la que girará el eje, así como el torque anteriormente calculado.

Al ser un sistema de baja velocidad constante de máximo 6 RPM, y si se sabe que:

\begin{center}
  $\omega = \frac{\theta}{t} $  
\end{center}
%Donde: \\
%$\theta$ son las vueltas en radianes \\
%$t$ es el tiempo en segundos
Se tiene entonces:
\begin{center}
    $\omega = \frac{(6 RPM)(2 \pi)}{60 s} rad/s$ \\
    $\omega = 0.6283 rad/s $ \\
    $\omega \approx 0.63 rad/s$
\end{center}

El cálculo de la potencia consumida se puede calcular con de la siguiente ecuación:

\begin{center}
 $P_{cons}=M_t * \omega$   
\end{center}

%Donde: \\

%$M_t$ es el Momento de torsión $\tau$ \\
%$\omega$ es la velocidad angular previamente calculada

Sustituyendo estos valores, se tiene que:
\begin{center}
    $P_{cons}=(40 Nm)(0.63 rad/s)$\\
    $P_{cons}=25.2 W$
\end{center}

